\chapter{La respiration et les échanges gazeux}

\noindent\textit{Respirer, c'est faire entrer de l'air riche en dioxygène jusque dans nos
poumons, et rejeter de l'air chargé en dioxyde de carbone. Ce chapitre étudie l'appareil
respiratoire, les mouvements qui l'animent, et les échanges gazeux qui s'y déroulent.}
\vspace{8pt}

\connecteBanner

\section{L'appareil respiratoire}

\begin{definition}[Trajet de l'air]
L'air inspiré traverse successivement : \textbf{nez ou bouche} $\to$ \textbf{pharynx}
$\to$ \textbf{larynx} $\to$ \textbf{trachée} $\to$ \textbf{bronches} $\to$
\textbf{bronchioles} $\to$ \textbf{alvéoles pulmonaires}, situées au bout des
bronchioles, à l'intérieur des \textbf{poumons}.
\end{definition}

\begin{illus}
\begin{tikzpicture}[scale=0.8]
\filldraw[onbo!20,draw=onbo,line width=1.1pt] (0,0) rectangle (3.2,0.9);
\node[font=\small] at (1.6,0.45) {nez / bouche};
\draw[->,onbmuted,line width=1.1pt] (3.2,0.45) -- (4.6,0.45);
\filldraw[onbgreen!20,draw=onbgreen,line width=1.1pt] (4.6,0) rectangle (6.8,0.9);
\node[font=\small] at (5.7,0.45) {trachée};
\draw[->,onbmuted,line width=1.1pt] (6.8,0.45) -- (8.2,0.45);
\filldraw[onbo2!20,draw=onbo2,line width=1.1pt] (8.2,0) rectangle (10.6,0.9);
\node[font=\small] at (9.4,0.45) {bronches};
\draw[->,onbmuted,line width=1.1pt] (10.6,0.45) -- (12,0.45);
\filldraw[onbink!15,draw=onbink,line width=1.1pt] (12,0) rectangle (16.7,0.9);
\node[font=\small] at (14.35,0.45) {alvéoles pulmonaires};
\end{tikzpicture}
\fig{Figure — Trajet de l'air : du nez aux alvéoles pulmonaires, lieu des échanges gazeux.}
\end{illus}

\section{Les mouvements respiratoires}

\begin{propriete}[Inspiration et expiration]
Lors de l'\textbf{inspiration}, le diaphragme et les muscles intercostaux se
\textbf{contractent} : le volume de la cage thoracique \textbf{augmente}, ce qui fait
entrer l'air dans les poumons. Lors de l'\textbf{expiration}, ces muscles se
\textbf{relâchent} : le volume thoracique \textbf{diminue}, ce qui chasse l'air hors des
poumons.
\end{propriete}

\begin{definition}[Fréquence respiratoire]
La \textbf{fréquence respiratoire} $f$ est le nombre de mouvements respiratoires par
minute (mvt/min) :
\[
f=\frac Nt,
\]
avec $N$ le nombre de mouvements comptés pendant une durée $t$ (en min).
\end{definition}

\begin{exemple}
Une personne effectue $N=18$ mouvements respiratoires en $t=1$ min : $f=18$ mvt/min.
\end{exemple}

\begin{remarque}
La fréquence respiratoire \textbf{augmente} pendant un effort physique : les muscles ont
alors besoin de davantage de dioxygène, et produisent davantage de dioxyde de carbone à
évacuer.
\end{remarque}

\section{Les échanges gazeux alvéolaires}

\begin{propriete}[Au niveau des alvéoles]
Chaque alvéole pulmonaire est entourée de nombreux \textbf{capillaires sanguins}. À
travers leur paroi très fine, le sang \textbf{capte le dioxygène} ($\mathrm{O_2}$) de
l'air alvéolaire et lui \textbf{cède le dioxyde de carbone} ($\mathrm{CO_2}$) qu'il
transportait. C'est un échange dans les \textbf{deux sens}, appelé \textbf{hématose}.
\end{propriete}

\begin{propriete}[Composition comparée de l'air]
\begin{center}
\small
\begin{tabular}{lcc}
\toprule
& \textbf{Air inspiré} & \textbf{Air expiré} \\
\midrule
Dioxygène $\mathrm{O_2}$ & $\approx21\,\%$ & $\approx16\,\%$ \\
Dioxyde de carbone $\mathrm{CO_2}$ & $\approx0{,}04\,\%$ & $\approx4\,\%$ \\
\bottomrule
\end{tabular}
\end{center}
L'air expiré est donc \textbf{appauvri en dioxygène} et \textbf{enrichi en dioxyde de
carbone} par rapport à l'air inspiré : la preuve que des échanges gazeux ont bien eu lieu
dans les poumons.
\end{propriete}

\section{Le débit ventilatoire}

\begin{propriete}[Débit ventilatoire]
Le \textbf{volume courant} $V_c$ est le volume d'air échangé à chaque mouvement
respiratoire (environ $0{,}5$ L au repos). Le \textbf{débit ventilatoire} $D$, volume
d'air total inspiré (ou expiré) par minute, vaut :
\[
D=V_c\times f.
\]
\end{propriete}

\begin{exemple}
Au repos, $V_c=0{,}5$ L et $f=16$ mvt/min : $D=0{,}5\times16=8$ L/min.
\end{exemple}

% ═══════════════════════════════════════════════════════════════
\section{L'essentiel à retenir}

\begin{synthese}
\textbf{Trajet de l'air.} Nez/bouche $\to$ pharynx $\to$ larynx $\to$ trachée $\to$
bronches $\to$ bronchioles $\to$ alvéoles.\\[4pt]
\textbf{Inspiration.} Muscles contractés, volume thoracique augmente. \textbf{Expiration.}
Muscles relâchés, volume diminue.\\[4pt]
\textbf{Fréquence respiratoire.} $f=\dfrac Nt$ (mvt/min).\\[4pt]
\textbf{Hématose.} Aux alvéoles : le sang capte $\mathrm{O_2}$ et cède $\mathrm{CO_2}$.\\[4pt]
\textbf{Débit ventilatoire.} $D=V_c\times f$.\\[4pt]
\textbf{Piège fréquent.} L'air expiré contient encore du dioxygène ($\approx16\,\%$) : il
n'est pas totalement épuisé, ce qui permet d'ailleurs le bouche-à-bouche.
\end{synthese}

% ═══════════════════════════════════════════════════════════════
\section{Méthodes \& savoir-faire}

\methode{Retracer le trajet de l'air dans l'appareil respiratoire}
Suivre l'ordre : nez/bouche, pharynx, larynx, trachée, bronches, bronchioles, alvéoles.
\begin{exemple}
L'air ne passe jamais directement du nez aux alvéoles sans traverser la trachée.
\end{exemple}

\methode{Calculer une fréquence respiratoire}
Appliquer $f=\dfrac Nt$, avec $t$ en minutes.
\begin{exemple}
$N=30$ mouvements en $t=1$ min : $f=30$ mvt/min.
\end{exemple}

\methode{Interpréter une comparaison air inspiré / air expiré}
Comparer les pourcentages de $\mathrm{O_2}$ et $\mathrm{CO_2}$ : une baisse de
$\mathrm{O_2}$ et une hausse de $\mathrm{CO_2}$ prouvent un échange gazeux dans les
poumons.
\begin{exemple}
Un air expiré à $16\,\%$ de $\mathrm{O_2}$ (contre $21\,\%$ inspiré) confirme une
consommation de dioxygène par l'organisme.
\end{exemple}

\methode{Calculer un débit ventilatoire}
Appliquer $D=V_c\times f$.
\begin{exemple}
$V_c=0{,}4$ L, $f=20$ mvt/min : $D=0{,}4\times20=8$ L/min.
\end{exemple}

% ═══════════════════════════════════════════════════════════════
\section{Exercices d'application}
\begin{multicols}{2}

\rubrique{Appareil respiratoire}
\exo{}\diff{1} Citer, dans l'ordre, les organes traversés par l'air entre le nez et les
alvéoles.

\exo{}\diff{2} Où se déroulent les échanges gazeux ?

\rubrique{Mouvements respiratoires}
\exo{}\diff{2} Que se passe-t-il au niveau des muscles lors de l'inspiration ?

\exo{}\diff{2} Que se passe-t-il au niveau du volume thoracique lors de l'expiration ?

\exo{}\diff{2} Une personne effectue $20$ mouvements respiratoires en $1$ min. Calculer
sa fréquence respiratoire.

\exo{}\diff{2} Une personne effectue $9$ mouvements en $0{,}5$ min. Calculer sa fréquence
respiratoire.

\rubrique{Échanges gazeux}
\exo{}\diff{2} Que capte le sang au niveau des alvéoles ? Que cède-t-il ?

\exo{}\diff{2} L'air inspiré contient environ $21\,\%$ de dioxygène ; l'air expiré en
contient environ $16\,\%$. Que peut-on en conclure ?

\exo{}\diff{1} Comment appelle-t-on l'échange gazeux qui a lieu au niveau des poumons ?

\rubrique{Débit ventilatoire}
\exo{}\diff{2} Calculer le débit ventilatoire pour $V_c=0{,}5$ L et $f=16$ mvt/min.

\exo{}\diff{2} Calculer le débit ventilatoire pour $V_c=0{,}6$ L et $f=15$ mvt/min.

% ═══════════════════════════════════════════════════════════════
\end{multicols}

\section{Exercices d'entraînement}
\begin{multicols}{2}

\rubrique{Appareil respiratoire}
\exo{}\diff{2} Pourquoi les alvéoles pulmonaires sont-elles très nombreuses et entourées
de nombreux capillaires ?

\exo{}\diff{3} Expliquer pourquoi une bronchite (inflammation des bronches) peut gêner la
respiration.

\rubrique{Mouvements respiratoires}
\exo{}\diff{3} Pourquoi la fréquence respiratoire augmente-t-elle pendant un effort
physique ?

\exo{}\diff{2} Une personne effectue $45$ mouvements en $3$ min. Calculer sa fréquence
respiratoire.

\exo{}\diff{3} Une personne au repos a une fréquence de $14$ mvt/min. Après un effort,
elle atteint $28$ mvt/min. Comparer les deux valeurs et proposer une explication.

\rubrique{Échanges gazeux}
\exo{}\diff{3} Pourquoi l'air expiré contient-il encore environ $16\,\%$ de dioxygène,
alors que l'organisme en consomme continuellement ?

\exo{}\diff{2} Pourquoi le taux de $\mathrm{CO_2}$ est-il plus élevé dans l'air expiré
que dans l'air inspiré ?

\exo{}\diff{3} Expliquer, à l'aide de la notion d'hématose, pourquoi un poisson hors de
l'eau ne peut pas respirer normalement (en une phrase, sans entrer dans le détail des
branchies).

\rubrique{Débit ventilatoire}
\exo{}\diff{2} Calculer le débit ventilatoire pour $V_c=0{,}5$ L et $f=25$ mvt/min.

\exo{}\diff{3} Un sportif a un débit ventilatoire $D=12$ L/min avec $V_c=0{,}6$ L.
Calculer sa fréquence respiratoire $f$.

\exo{}\diff{3} Expliquer pourquoi le débit ventilatoire augmente pendant un effort
physique, en citant les deux grandeurs qui peuvent varier.

% ═══════════════════════════════════════════════════════════════
\end{multicols}

\section{Sujets type BEPC}

\sujetbac[Trajet de l'air et échanges gazeux]
\begin{enumerate}[label=\arabic*)]
\item Cite, dans l'ordre, les organes traversés par l'air entre le nez et les alvéoles.
\item Où se déroulent les échanges gazeux respiratoires ?
\item Le tableau ci-dessous donne la composition de l'air inspiré et de l'air expiré.
\begin{center}
\small
\begin{tabular}{lcc}
\toprule
& Air inspiré & Air expiré \\
\midrule
$\mathrm{O_2}$ & $21\,\%$ & $16\,\%$ \\
$\mathrm{CO_2}$ & $0{,}04\,\%$ & $4\,\%$ \\
\bottomrule
\end{tabular}
\end{center}
Que peut-on en conclure sur les échanges gazeux dans les poumons ?
\item Comment appelle-t-on cet échange gazeux ?
\end{enumerate}

\sujetbac[Fréquence et débit ventilatoire]
\begin{enumerate}[label=\arabic*)]
\item Donne la relation permettant de calculer la fréquence respiratoire $f$.
\item Un élève compte $16$ mouvements respiratoires en $1$ min au repos. Calcule sa
fréquence respiratoire.
\item Après une course, il compte $32$ mouvements en $1$ min. Calcule la nouvelle
fréquence, et compare-la à celle du repos.
\item Sachant que son volume courant est $V_c=0{,}5$ L, calcule son débit ventilatoire
au repos puis après l'effort.
\end{enumerate}

\sujetbac[Mouvements respiratoires]
\begin{enumerate}[label=\arabic*)]
\item Que se passe-t-il au niveau du diaphragme et des muscles intercostaux lors de
l'inspiration ?
\item Quelle conséquence cela a-t-il sur le volume de la cage thoracique ?
\item Décris ce qui se passe lors de l'expiration.
\item Pourquoi la respiration s'accélère-t-elle pendant un effort physique ?
\end{enumerate}

% ═══════════════════════════════════════════════════════════════
\section{Corrigés}
\begin{multicols}{2}

\corrige{de l'exercice 1}
Nez/bouche, pharynx, larynx, trachée, bronches, bronchioles, alvéoles.

\corrige{de l'exercice 2}
Au niveau des alvéoles pulmonaires.

\corrige{de l'exercice 3}
Ils se contractent.

\corrige{de l'exercice 4}
Il diminue.

\corrige{de l'exercice 5}
$f=\dfrac{20}1=20$ mvt/min.

\corrige{de l'exercice 6}
$f=\dfrac9{0{,}5}=18$ mvt/min.

\corrige{de l'exercice 7}
Il capte le dioxygène et cède le dioxyde de carbone.

\corrige{de l'exercice 8}
L'organisme consomme du dioxygène : une partie de l'$\mathrm{O_2}$ inspiré est captée par
le sang au niveau des poumons.

\corrige{de l'exercice 9}
L'hématose.

\corrige{de l'exercice 10}
$D=0{,}5\times16=8$ L/min.

\corrige{de l'exercice 11}
$D=0{,}6\times15=9$ L/min.

\corrige{du sujet 1}
1) Nez/bouche, pharynx, larynx, trachée, bronches, bronchioles, alvéoles.\\
2) Au niveau des alvéoles pulmonaires.\\
3) Le sang a capté du dioxygène (baisse de $21\,\%$ à $16\,\%$) et a cédé du dioxyde de
carbone (hausse de $0{,}04\,\%$ à $4\,\%$).\\
4) L'hématose.

\corrige{du sujet 2}
1) $f=\dfrac Nt$.\\
2) $f=\dfrac{16}1=16$ mvt/min.\\
3) $f=\dfrac{32}1=32$ mvt/min ; elle a doublé par rapport au repos.\\
4) Au repos : $D=0{,}5\times16=8$ L/min. Après l'effort : $D=0{,}5\times32=16$ L/min.

\corrige{du sujet 3}
1) Ils se contractent.\\
2) Le volume de la cage thoracique augmente, ce qui fait entrer l'air.\\
3) Les muscles se relâchent, le volume thoracique diminue, l'air est chassé hors des
poumons.\\
4) Car les muscles en activité ont besoin de davantage de dioxygène et produisent
davantage de dioxyde de carbone à évacuer.

\end{multicols}
\exercicesBanner
